\documentclass[12pt]{article}

\usepackage[spanish]{babel}
\usepackage[utf8]{inputenc}
\usepackage{geometry}
\usepackage{minted}

\geometry{margin=2.5cm}

% TÍTULO CORRECTO
\begin{titlepage}
    \centering
    
\title{Laboratorio 1 \\ Fundamentos de Informática \\ Universidad Nacional}
\author{David Antonio Zúñiga Larcón}
\date{26 de marzo}

\maketitle

\end{titlepage}  

\begin{document}
% CONTENIDO

\section{Intro}


\section{Código en C++}

Ejemplo de uso con \verb|cout|:

\begin{minted}{cpp}
#include <iostream>
using namespace std;

int main() {
    cout << "Hola mundo" << endl;
    return 0;
}
\end{minted}

\section{}


\section{}



\end{document}

 \begin{figure}[h!] 
 para imagenes: \centering \includegraphics [scale=0.6]{Imganen1.png} % Nombre de la imagen para la segunda ley. \caption{Representación gráfica de la segunda ley de Kepler, (Serway y Jewett, 2008, p. 368).} \label{fig:kepler2} \end{figure verb 
 
 \verb|cout << "Hola";| \verb!cout << "Hola";! 
 
 \setminted{
linenos,
frame=single,
framesep=0mm,
rulecolor=\color{blue},
bgcolor=gray!10,
fontsize=\small
... el resto de opciones que se desean cargar
4. Ejemplo con opciones de formato
La librería minted permite personalizar la apariencia del bloque de código; entre las
opciones mas frecuentes están:
Opciones visuales básicas del bloque
• fontsize=\small: cambia el tamaño de letra.
• linenos: muestra números de línea.
• frame=single: pone un recuadro.
• framesep=...: controla el espacio entre el borde y el código.
• baselinestretch=...: cambia el espaciado vertical entre líneas.
• bgcolor=...: pone color de fondo.
• tabsize=2: fija el tamaño visual de las tabulaciones.
• breaklines: parte líneas largas automáticamente.
• breakanywhere: permite cortes en más posiciones.
• breaksymbolleft, breaksymbolright: cambian el símbolo que aparece cuando
una línea se parte.

Opciones de las líneas de numeración
• linenos: activa numeración.
• firstnumber=1: fija el número inicial.
• stepnumber=2: numera cada 2 líneas.
• numbersep=5pt: separa los números del código.
• numberblanklines=true/false: decide si líneas en blanco se numeran.
• resetmargins: ajusta márgenes cuando hay numeración.
Opciones de márgenes, sangría y caja
• xleftmargin=...
• xrightmargin=...
• gobble=2: elimina una cantidad fija de espacios al inicio de cada línea.
• autogobble: detecta y elimina sangría común automáticamente.
• samepage: intenta no partir el bloque entre páginas.
• frameleftline, framerightline, frametopline,framebottomline: controlan
lados específicos del marco.
Opciones del resaltado de sintaxis
• style=default, style=monokai, etc.: cambia el estilo de colores.
• {cpp}, {python}, {java}, etc.: el lenguaje se indica en el entorno.
• encoding=utf8: define codificación cuando aplica.
• texcomments: interpreta comentarios para permitir comandos TeX en ciertos
casos.
• mathescape: permite escribir matemáticas dentro del código usando ....
• escapeinside=|| o similar: permite “salir” del código y escribir LaTeX dentro
del bloque.
• highlightlines={2,4-6}:
Opciones para inclusión de archivos
Al usar \inputminted se puede agregar opciones para leer solo partes de un archivo.
• firstline=...
• lastline=...
• firstnumber=...
• stripall
• stripnl
• gobble=...
Estas opciones son muy útiles cuando el código está en archivos externos.
Para definir un entorno minted con opciones se usa la siguiente estructura
\begin{minted}[linenos, frame=single, bgcolor=gray!10, resto de Opciones]{cpp}
// codigo c++
\end{minted}